% Adaptación de «Red, residuales y pérdidas» de arquitectura_pinn.html.
\begin{table}[p]
\centering
\caption{Notación de la arquitectura y del entrenamiento de la PINN condicional.}
\label{met:tabla_pinn_notacion}
\begingroup
\small
\setlength{\tabcolsep}{4pt}
\renewcommand{\arraystretch}{1.12}
\begin{tabularx}{\textwidth}{@{}>{\raggedright\arraybackslash}p{0.20\textwidth}>{\raggedright\arraybackslash}X>{\raggedright\arraybackslash}p{0.28\textwidth}@{}}
\toprule
Símbolo & Significado & Convención \\
\midrule
$\eta;\ \mathcal N_\eta$ & Pesos y sesgos neuronales; red que produce el incremento normalizado. & Se ajustan mediante optimización. \\
$\mathbf z;\ n_z$ & Entrada concatenada de la red en la Figura~\ref{met:fig_pinn_arquitectura} y su dimensión. & $n_z=1+4+\dim\widetilde{\mathbf u}+\dim\widetilde{\mathbf d}$. \\
$L$ (capas) & Profundidad de la red en el esquema; se ilustran la primera y la última capa oculta. & Profundidad y anchuras sujetas a selección. \\
Tilde $(\widetilde{\ \ })$ & Entrada normalizada mediante transformaciones estimadas en desarrollo. & La transformación se conserva en validación y prueba. \\
$\mathbf s_x;\ s_q$ & Escalas físicas de los estados y de la componente $q$. & K para temperaturas; kg/kg para $W$. \\
$\odot;\ \delta\mathbf x$ & Producto por componente e incremento del estado en unidades físicas. & $\delta\mathbf x=(s/\Delta t_c)\mathcal N_\eta(\mathbf z)\odot\mathbf s_x$. \\
$\boldsymbol\theta;\ \mathbf f$ & Parámetros físicos y derivadas de estado obtenidas de los balances. & $\mathbf f=[\dot T_c,\dot T_b,\dot T_a,\dot W]^{\mathsf T}$. \\
$q;\ r_q$ & Índice de componente y residual de su ecuación diferencial. & $q\in\{T_c,T_b,T_a,W\}$; residual en estado/s. \\
$s_\ell;\ N_f$ & Tiempo del punto de colocación $\ell$ y número de puntos. & Dominio físicamente admitido. \\
$\mathcal H_j;\ y_j$ & Operador de observación del sensor $j$ y observación disponible. & Magnitud, posición y soporte compatibles. \\
$\mathcal L_D,\mathcal L_F,\mathcal L_R$ & Pérdidas de datos, física y restricciones. & Condición inicial impuesta por estructura. \\
$\lambda_D,\lambda_F,\lambda_R$ & Pesos de las tres contribuciones de pérdida. & Se documentan con las escalas y la selección en desarrollo. \\
\bottomrule
\end{tabularx}
\endgroup
\par\smallskip
\begin{minipage}{\textwidth}\footnotesize
$\mathbf z$ y $L$ se emplean aquí como notación local del esquema de arquitectura: entrada normalizada y profundidad de la red. Los parámetros neuronales $\eta$ se distinguen de los parámetros físicos $\boldsymbol\theta$; estos últimos pueden fijarse por identificación o estimarse en desarrollo con cotas e identificabilidad.
\end{minipage}
\end{table}
